% Sources:
% - docs/paper/threat-model-draft.md
% - docs/paper/problem-statement-draft.md
% - docs/paper/contribution-evidence-map.md
% - docs/rq-v2/rq-v2-threat-model.md
% - docs/rq-v2/g2-admission-semantics.md
% - docs/rq-v2/rq-v2-complete-argument-map.md

\subsection{Party-Level Batch Identity Objective}

Section~2 identified K-Waay's stated distinct-party condition and separated
party identity from message, occurrence, slot, and batch coordinates. We now
formalize the same-batch relation used as the analysis objective. Let
\(B=(E_1,\ldots,E_n)\) be an ordered batch and let
\(\mathsf{party}(E_i)\) denote the modeled principal associated with entry
\(E_i\). We define
\begin{equation}
  \distinctparty(B)
  \mathrel{\triangleq}
  \forall\,1 \leq i < j \leq n,\;
  \mathsf{party}(E_i)\neq\mathsf{party}(E_j).
  \label{eq:distinct-party-objective}
\end{equation}
Every pair of distinct positions in one admitted batch must therefore project
to distinct modeled principals. This is a relation over party coordinates; it
is not obtained by comparing message contents, occurrence identifiers, or slot
indices.

The objective is local to one batch. It does not require a party to occur only
once throughout the protocol, across batches, or across sessions, and it does
not impose global uniqueness on messages or sender occurrences. We describe a
composition that fails Equation~\eqref{eq:distinct-party-objective} only as
violating the party-level batch invariant. Such a violation does not by itself
imply malformed cryptographic material, authentication failure, or a break of
an underlying primitive.

The invariant is therefore a structural predicate on the composition admitted
as one batch, rather than a cryptographic authenticity predicate for an
individual entry. It states which party projections may coexist, while leaving
the mechanism that obtains and checks those projections outside the objective
itself.

Equation~\eqref{eq:distinct-party-objective} is stated for arbitrary finite
\(n\) to define the intended relation independently of the model bound. The
machine-checked analysis fixes \(n=2\), where the objective reduces to
\(\mathsf{party}(E_1)\neq\mathsf{party}(E_2)\). We claim no full-protocol proof
for arbitrary batch sizes. At this point we also do not assume that violating
the invariant has an observable receiver-side consequence; Sections~4 and~5
define and evaluate that question.

\subsection{Batch-Composition Adversary}

We consider a symbolic batch-composition adversary. Control over the selection
and placement of entries is an analytical capability assumption at the modeled
composition boundary, not a claim that a deployed network attacker controls
the corresponding batching component. The adversary acts at the interface between
the availability of sender-associated candidate entries and the decision to
admit their composition as one batch. Candidate tuples are produced by the
modeled sender lifecycle and exposed on the symbolic network. The adversary may
retain and replay an exposed tuple, select among exposed tuples, arrange
selected tuples in either batch position, and present the resulting
composition to the admission step. These capabilities concern input
composition; they do not grant control over the creation of honest-party
state.

The threat model distinguishes two forms of repeated input. First, the
adversary may present the same exposed tuple twice, so that \(E_1=E_2\).
Second, it may select two distinct sender-origin tuples
\begin{equation}
  E_1=(A,\mathit{sid}_1,m_1),
  \qquad
  E_2=(A,\mathit{sid}_2,m_2),
  \qquad
  E_1\neq E_2,
  \label{eq:repeated-party-capability}
\end{equation}
that nevertheless share the party projection \(A\). In the second case, the
tuples differ in both occurrence and message under the prototype's fresh-value
assumption. The capability is therefore
one of batch composition rather than identity forgery: it does not require the
adversary to impersonate or invent a second modeled principal.

For claims about an accepted component, sender--receiver matching is defined
over the complete modeled origin tuple \((A,\mathit{sid},m)\), not over \(A\)
alone.
Although an adversary may syntactically recombine values learned from the
network, a recombined tuple does not have a matching sender origin unless that
exact tuple was produced by the modeled sender lifecycle. Party equality by
itself therefore supplies no origin evidence. Section~4 gives the concrete
event and state encoding of this relation.

The analyzed composition does not rely on key recovery, a break of the
underlying KEM or signature mechanisms, or forgery of cryptographic material
abstracted away by the model. These exclusions bound the abstraction rather
than proving the omitted mechanisms secure. The model also does not determine
how a concrete implementation exposes, buffers, or schedules corresponding
wire-format inputs. Whether an available composition passes an admission
predicate or leads to later receiver events is a verification question, not an
assumption of this threat model. Section~4 compares admission predicates
expressed over different identity coordinates.

\subsection{Analysis Boundary}

The direct formal object contains the sender lifecycle needed to produce
origin tuples, symbolic network composition, batch admission or rejection,
the sender-origin relation, and component-level \receiveraccept events. The
origin relation used in correspondence and injectivity properties matches the
full \((A,\mathit{sid},m)\) tuple. Each \receiveraccept event is a modeling
event local to our admission abstraction and should not be read as a
protocol-level acceptance event defined by K-Waay. Introducing it assumes no
verification outcome.

Direct machine-checked evidence ends at \receiveraccept. The model omits
K-Waay's cryptographic computations, compromise behavior, downstream key
interfaces, ratcheting state, and application consumers. In particular, it
contains no formal \textsc{Key} or \textsc{Test} objects and does not directly
establish confidentiality, key indistinguishability, or full-protocol
authentication failure.

Consequences for upper-layer state installation would require an additional
composition assumption connecting each modeled acceptance occurrence to that
consumer and are not asserted here. The analysis likewise makes no claim about
a deployed implementation, an exploitable wire-level vulnerability, or the
validity of K-Waay's original cryptographic theorem. Under this boundary, we
study the research questions stated earlier; Section~4 defines their symbolic
realization.
